\section{Related Work}
\label{sec:related}

\paragraph*{Interruption and resumption.}
The resumption lag construct, the interval between the end of an interruption
and the resumption of productive work on the primary task, is the standard
dependent measure in this
literature~\citep{brandt2011resumption,ferreiro2014fragmentation}. Its
magnitude depends on the depth of the suspended state and on whether the
interruption arrives at a subtask boundary. Work on peripheral and ambient
displays argues that the cost is driven by attention shifts rather than by
information volume~\citep{okonjo2016peripheral}, which is the mechanism our
conditions are designed to separate.

\paragraph*{Developer interruption in the field.}
Field studies of programmers report interruption rates of several per hour and
recovery times in the range of ten to twenty minutes for deep
interruptions~\citep{ferreiro2014fragmentation,steinbach2019telemetry}. Those
figures come from telemetry over long windows and cannot isolate the cost of
one notification design, because a field deployment cannot hold task difficulty
constant. Our contribution is complementary: a controlled setting that trades
ecological validity for the ability to attribute a difference to modality.

\paragraph*{Notification design.}
Deferring notifications to breakpoints reduces their
cost~\citep{ravichandran2018deferral}, and batching them reduces it
further~\citep{nkemdirim2020batching}, but both of these delay the information.
For a build result the delay is the thing under test, since a developer waiting
on a compile has a reason to want the result promptly. Our batched condition
quantifies that trade directly.

\paragraph*{Compile feedback specifically.}
Studies of continuous compilation and incremental type checking focus on
latency and correctness rather than on how the result is
presented~\citep{lindqvist2017continuous}. To our knowledge the presentation
question has not been tested experimentally in this setting, which is the gap
this paper fills.
